
\section{Hilbert folding: from tick time to holographic screen locality}
\label{sec:hilbert_folding}

\subsection{Addressing at finite resolution}
\label{sec:finite_resolution_addressing}

Fix $d\ge 2$.
At resolution level $n$, the screen lattice $\Sigma_n$ in \eqref{eq:screen_lattice} has $2^{dn}$ sites.
By Axiom~\ref{ax:H1}, there exists a bijective address map $H_n$ with one-step locality \eqref{eq:hilbert_locality}.
We interpret $H_n$ as a \emph{compiler-visible} and \emph{auditable} specification of how scan ticks visit screen degrees of freedom.

Because $H_n$ is a bijection, it admits an inverse
\[
H_n^{-1}:\Sigma_n\to\{0,1,\dots,2^{dn}-1\},
\]
and therefore supports two primitive operations:
\begin{itemize}
\item \textbf{forward addressing} (tick $\to$ site): $t\mapsto H_n(t)$;
\item \textbf{reverse addressing} (site $\to$ tick): $x\mapsto H_n^{-1}(x)$.
\end{itemize}
This invertibility is essential: locality is not only a geometric notion but a compilable addressing relation.

\subsection{Locality as a verifiable adjacency constraint}
\label{sec:hilbert_locality_property}

The defining local constraint \eqref{eq:hilbert_locality} implies that the scan visits screen sites along a nearest-neighbor path:
\[
H_n(0)\sim H_n(1)\sim \cdots \sim H_n(2^{dn}-1),
\]
where $\sim$ denotes $\ell^1$ nearest-neighbor adjacency on the grid.
Thus, any scan-layer update that couples consecutive ticks corresponds to a screen-local propagation.

Conversely, implementing a prescribed screen-local interaction (nearest neighbors on $\Sigma_n$) on a one-dimensional scan substrate typically requires \emph{routing}:
degrees of freedom must be swapped/moved in scan order so that neighboring sites become neighboring ticks.
This is the source of computational overhead $\kappa(x)$ and, ultimately, computational lapse (Section~\ref{sec:gravity}).

\subsection{A constructive view of ``space''}
\label{sec:space_as_addressing}

Within this framework, ``space'' is defined at each resolution level $n$ as the graph $(\Sigma_n,\sim)$, where $\sim$ is a chosen nearest-neighbor relation (e.g.\ $\ell^1$).
The scan is not embedded \emph{into} a pre-existing manifold; rather, the manifold is an \emph{effective} description of how the address-induced neighborhood structure behaves under coarse graining.

Two consequences are immediate:
\begin{itemize}
\item \textbf{Locality is protocol-relative.}
The same underlying scan dynamics can induce different effective spatial neighborhoods under different admissible address maps (different compilation conventions).
\item \textbf{Continuum geometry is an emergent coarse-grained object.}
The continuous Hilbert curve $H:[0,1]\to[0,1]^d$ provides a classical scaling limit of $H_n$ under appropriate normalization, but operational claims are always made at finite $n$ and finite readout resolution.
\end{itemize}

\paragraph{A quantitative locality fact (H\"older regularity).}
For the continuous Hilbert curve $H:[0,1]\to[0,1]^d$, one has H\"older continuity with exponent $1/d$:
there exists a constant $C_d$ such that
\[
\norm{H(t)-H(s)}_\infty \le C_d\,|t-s|^{1/d}
\qquad (s,t\in[0,1]),
\]
see, e.g., \cite{Sagan1994}.
This provides a quantitative backbone for the ``finite resolution'' viewpoint: small parameter increments induce small spatial displacements after coarse graining, with a deterministic scale law.

\subsection{Algorithmic realization (2D) and locality check}
\label{sec:hilbert_algorithm}

For $d=2$, a standard bitwise construction provides a concrete $H_n$ and is easy to verify by brute force.
Appendix~\ref{app:hilbert_addressing} records a pure-Python implementation and checks \eqref{eq:hilbert_locality} for orders up to $n=8$ in milliseconds.
This implementation is included only as an auditable reference construction consistent with the standard literature \cite{Sagan1994}.


